\prob{1}{fluxo1} São dados um grafo $G=(V, E)$ direcionado, dois vértices $s$ (source) e $t$ (sink) e uma taxa máxima de fluxo (capacidade) pra cada aresta. Determinar o fluxo máximo.

\textbf{Soluções:}
\begin{enumerate}
    \item Ford-Fulkerson: $O(EF)$, onde $F$ é o fluxo máximo.
    \item Edmonds-Karp: $O(V E^2)$. Ou use Dinic:
    \item Dinic: $O(V^2 E)$
\end{enumerate}

\textbf{Edmonds-Karp:}

(Use uma dfs para virar Ford-Fulkerson)

\begin{lstlisting}[language=C++, breaklines=true]
int n;
vector<vector<int>> capacity;
vector<vector<int>> adj;

int bfs(int s, int t, vector<int>& parent) {
    fill(parent.begin(), parent.end(), -1);
    parent[s] = -2;
    queue<pair<int, int>> q;
    q.push({s, INF});

    while (!q.empty()) {
        int cur = q.front().first;
        int flow = q.front().second;
        q.pop();

        for (int next : adj[cur]) {
            if (parent[next] == -1 && capacity[cur][next]) {
                parent[next] = cur;
                int new_flow = min(flow, capacity[cur][next]);
                if (next == t)
                    return new_flow;
                q.push({next, new_flow});
            }
        }
    }

    return 0;
}

int maxflow(int s, int t) {
    int flow = 0;
    vector<int> parent(n);
    int new_flow;

    while (new_flow = bfs(s, t, parent)) {
        flow += new_flow;
        int cur = t;
        while (cur != s) {
            int prev = parent[cur];
            capacity[prev][cur] -= new_flow;
            capacity[cur][prev] += new_flow;
            cur = prev;
        }
    }

    return flow;
}
\end{lstlisting}

\textbf{Dinic:}

\begin{lstlisting}[language=C++, breaklines=true]
struct FlowEdge {
    int v, u;
    long long cap, flow = 0;
    FlowEdge(int v, int u, long long cap) : v(v), u(u), cap(cap) {}
};

struct Dinic {
    const long long flow_inf = 1e18;
    vector<FlowEdge> edges;
    vector<vector<int>> adj;
    int n, m = 0;
    int s, t;
    vector<int> level, ptr;
    queue<int> q;

    Dinic(int n, int s, int t) : n(n), s(s), t(t) {
        adj.resize(n);
        level.resize(n);
        ptr.resize(n);
    }

    void add_edge(int v, int u, long long cap) {
        edges.emplace_back(v, u, cap);
        edges.emplace_back(u, v, 0);
        adj[v].push_back(m);
        adj[u].push_back(m + 1);
        m += 2;
    }

    bool bfs() {
        while (!q.empty()) {
            int v = q.front();
            q.pop();
            for (int id : adj[v]) {
                if (edges[id].cap - edges[id].flow < 1)
                    continue;
                if (level[edges[id].u] != -1)
                    continue;
                level[edges[id].u] = level[v] + 1;
                q.push(edges[id].u);
            }
        }
        return level[t] != -1;
    }

    long long dfs(int v, long long pushed) {
        if (pushed == 0)
            return 0;
        if (v == t)
            return pushed;
        for (int& cid = ptr[v]; cid < (int)adj[v].size(); cid++) {
            int id = adj[v][cid];
            int u = edges[id].u;
            if (level[v] + 1 != level[u] || edges[id].cap - edges[id].flow < 1)
                continue;
            long long tr = dfs(u, min(pushed, edges[id].cap - edges[id].flow));
            if (tr == 0)
                continue;
            edges[id].flow += tr;
            edges[id ^ 1].flow -= tr;
            return tr;
        }
        return 0;
    }

    long long flow() {
        long long f = 0;
        while (true) {
            fill(level.begin(), level.end(), -1);
            level[s] = 0;
            q.push(s);
            if (!bfs())
                break;
            fill(ptr.begin(), ptr.end(), 0);
            while (long long pushed = dfs(s, flow_inf)) {
                f += pushed;
            }
        }
        return f;
    }
};
\end{lstlisting}

\prob{2}{fluxo2} $s-t-\textbf{cut}$ é uma partição dos vértices em um conjunto que contém $s$ e outro $t$. A capacidade de um cut é a soma das capacidades das arestas indo de um conjunto pra outro. É conhecido que a capacidade do max flow é igual a do min cut. Ache esse min-cut.

\textbf{Solução:} Faça o max-flow. Particione em: conjunto de vértices alcançados por $s$ no grafo residual e o resto.

\prob{3}{fluxo3} Fluxo com demandas. Para cada $e \in E$, queremos achar uma função $f$ tal que $d(e) \leq f(e) \leq c(e)$ em que $d, c$ são funções dadas.

\textbf{Solução:} Adicione uma source $s'$ e um sink $t'$, uma aresta de $s'$ para todos os outros vértices, uma nova aresta de cada vértice para o sink $t'$ e uma aresta de $t$ para $s$. Defina uma nova função capacidade tal que

\begin{enumerate}
    \item $c'((s', v)) = \sum_{u \in V} d((u, v))$ para cada aresta $(s', v)$.
    \item $c'((v, t')) = \sum_{w \in V} d((v, w))$ para cada aresta $(v, t')$
    \item $c'((u, v)) = c((u, v)) - d((u, v))$ para cada aresta $(u, v)$ na rede antigo.
    \item $c'((t, s)) = \infty$
\end{enumerate}

Se essa nova rede tiver um fluxo saturante (em que cada aresta saindo de $s'$ está completamente preenchida, que dá no mesmo que as indo pra $t'$ estarem), então a rede com demandas tem um fluxo válido. Basta usar algum algoritmo de max-flow para achar o saturating-flow. 

\prob{4}{fluxo4} Idêntico ao \ref{prob:fluxo3}, mas queremos o fluxo mínimo.

\textbf{Solução:} Ao invés de usar $c'((t, s)) = \infty$, faça uma busca binária em $c'((t, s))$.

\prob{5}{fluxo5} (min-cost flow) Para cada aresta, são dados a capacidade e o custo por unidade de fluxo. Encontre um fluxo de $K$ e entre todos esses fluxos encontre o com menor custo.

\textbf{Solução:}

\begin{lstlisting}[language=C++, breaklines=true]
struct Edge
{
    int from, to, capacity, cost;
};

vector<vector<int>> adj, cost, capacity;

const int INF = 1e9;

void shortest_paths(int n, int v0, vector<int>& d, vector<int>& p) {
    d.assign(n, INF);
    d[v0] = 0;
    vector<bool> inq(n, false);
    queue<int> q;
    q.push(v0);
    p.assign(n, -1);

    while (!q.empty()) {
        int u = q.front();
        q.pop();
        inq[u] = false;
        for (int v : adj[u]) {
            if (capacity[u][v] > 0 && d[v] > d[u] + cost[u][v]) {
                d[v] = d[u] + cost[u][v];
                p[v] = u;
                if (!inq[v]) {
                    inq[v] = true;
                    q.push(v);
                }
            }
        }
    }
}

int min_cost_flow(int N, vector<Edge> edges, int K, int s, int t) {
    adj.assign(N, vector<int>());
    cost.assign(N, vector<int>(N, 0));
    capacity.assign(N, vector<int>(N, 0));
    for (Edge e : edges) {
        adj[e.from].push_back(e.to);
        adj[e.to].push_back(e.from);
        cost[e.from][e.to] = e.cost;
        cost[e.to][e.from] = -e.cost;
        capacity[e.from][e.to] = e.capacity;
    }

    int flow = 0;
    int cost = 0;
    vector<int> d, p;
    while (flow < K) {
        shortest_paths(N, s, d, p);
        if (d[t] == INF)
            break;

        // find max flow on that path
        int f = K - flow;
        int cur = t;
        while (cur != s) {
            f = min(f, capacity[p[cur]][cur]);
            cur = p[cur];
        }

        // apply flow
        flow += f;
        cost += f * d[t];
        cur = t;
        while (cur != s) {
            capacity[p[cur]][cur] -= f;
            capacity[cur][p[cur]] += f;
            cur = p[cur];
        }
    }

    if (flow < K)
        return -1;
    else
        return cost;
}
\end{lstlisting}

\prob{6}{fluxo6} Algoritmo húngaro com fluxo. Selecionar $n$ elementos de uma matriz quadrada $A$ de ordem $n$, com nenhuma fileira em comum, de modo a maximizar sua soma.

\textbf{Solução:}

\begin{lstlisting}[language=C++, breaklines=true]
const int INF = 1000 * 1000 * 1000;

vector<int> assignment(vector<vector<int>> a) {
    int n = a.size();
    int m = n * 2 + 2;
    vector<vector<int>> f(m, vector<int>(m));
    int s = m - 2, t = m - 1;
    int cost = 0;
    while (true) {
        vector<int> dist(m, INF);
        vector<int> p(m);
        vector<bool> inq(m, false);
        queue<int> q;
        dist[s] = 0;
        p[s] = -1;
        q.push(s);
        while (!q.empty()) {
            int v = q.front();
            q.pop();
            inq[v] = false;
            if (v == s) {
                for (int i = 0; i < n; ++i) {
                    if (f[s][i] == 0) {
                        dist[i] = 0;
                        p[i] = s;
                        inq[i] = true;
                        q.push(i);
                    }
                }
            } else {
                if (v < n) {
                    for (int j = n; j < n + n; ++j) {
                        if (f[v][j] < 1 && dist[j] > dist[v] + a[v][j - n]) {
                            dist[j] = dist[v] + a[v][j - n];
                            p[j] = v;
                            if (!inq[j]) {
                                q.push(j);
                                inq[j] = true;
                            }
                        }
                    }
                } else {
                    for (int j = 0; j < n; ++j) {
                        if (f[v][j] < 0 && dist[j] > dist[v] - a[j][v - n]) {
                            dist[j] = dist[v] - a[j][v - n];
                            p[j] = v;
                            if (!inq[j]) {
                                q.push(j);
                                inq[j] = true;
                            }
                        }
                    }
                }
            }
        }

        int curcost = INF;
        for (int i = n; i < n + n; ++i) {
            if (f[i][t] == 0 && dist[i] < curcost) {
                curcost = dist[i];
                p[t] = i;
            }
        }
        if (curcost == INF)
            break;
        cost += curcost;
        for (int cur = t; cur != -1; cur = p[cur]) {
            int prev = p[cur];
            if (prev != -1)
                f[cur][prev] = -(f[prev][cur] = 1);
        }
    }

    vector<int> answer(n);
    for (int i = 0; i < n; ++i) {
        for (int j = 0; j < n; ++j) {
            if (f[i][j + n] == 1)
                answer[i] = j;
        }
    }
    return answer;
}
\end{lstlisting}